% Audited prompt documentation for the relational-reasoning population game.
% Kept as a separate appendix fragment so that the prompt can be revised and
% reviewed independently of the main methods text.

\definecolor{relpromptaccent}{HTML}{E9A126}
\definecolor{relpromptprivate}{HTML}{157F4A}
\definecolor{relpromptsocial}{HTML}{1F5FA9}
\definecolor{relpromptrule}{HTML}{C8C8C8}

\newcommand{\relvariable}[1]{\textcolor{relpromptaccent}{#1}}
\newcommand{\relprivate}[1]{\textcolor{relpromptprivate}{#1}}
\newcommand{\relsocial}[1]{\textcolor{relpromptsocial}{#1}}

\newtcolorbox{relpromptbox}[1]{%
  enhanced,
  breakable,
  colback=white,
  colframe=relpromptrule,
  boxrule=0.6pt,
  arc=2mm,
  left=7pt,
  right=7pt,
  top=6pt,
  bottom=6pt,
  fonttitle=\bfseries\small\color{relpromptaccent},
  colbacktitle=black!5,
  coltitle=relpromptaccent,
  boxed title style={boxrule=0pt,arc=0pt},
  title={#1}}

\section{Relational-reasoning prompt}
\label{app:relational-prompt}

The population is driven by one versioned prompt family,
\texttt{relational\_public\_ballot} version~1. Each call is compiled into two
chat messages. The system message establishes the agent's role, privacy
boundary, task, call-specific option mapping, and currently known facts. The
user message supplies the current vote and visible social source, then states
the response contract. The model therefore receives only the state available
to the focal agent at that microscopic update; it is never shown the global
vote vector, other agents' private facts, the ground-truth answer, or the
controller's sensor sample.

The display below reproduces one complete prompt from the audit of the minimal
Study~04 inspection episode ($N=24$, $q_c=12$, $b=12$). It is the first focal
update of population round~1 for Agent~17. Typography is annotation only:
\relvariable{orange} marks values instantiated anew for the call,
\relprivate{green} marks the focal agent's private state, and
\relsocial{blue} marks the visible social input. None of these colours or the
``SYSTEM MESSAGE''/``USER MESSAGE'' labels are sent to the model.

\begin{relpromptbox}{Audited focal-update prompt}
\small\raggedright
\setlength{\parindent}{0pt}
\textbf{SYSTEM MESSAGE}

You are \relvariable{Agent 17}, one participant in a group reasoning problem.

Your goal is to identify the correct answer.

Participant identities persist across interactions.

Different participants know different facts about this problem. No single
participant is guaranteed to know everything needed to answer, and some
participants may have objectives that differ from yours. You do not know
which participants, if any, have different objectives.

Make your own decision, using:
\begin{itemize}[leftmargin=1.3em,itemsep=0pt,topsep=2pt]
  \item the facts you currently know;
  \item your current position;
  \item the public positions of the participants shown below.
\end{itemize}

You see only the participants shown in this interaction, and of each one you
see their vote and the fact they chose to share, if any.

Other participants will see the same of you: your vote and any fact you choose
to share, and nothing else. Your reason is your own record: it is not shown to
anyone.

\textbf{QUESTION}

\relvariable{Where is Davi relative to Hira?}

The available answers are:
\begin{itemize}[leftmargin=1.3em,itemsep=0pt,topsep=2pt]
  \item \relvariable{A) WEST}
  \item \relvariable{B) EAST}
  \item \relvariable{C) SOUTHWEST}
\end{itemize}

Exactly one of these answers is correct. Vote by its letter.

\textbf{YOUR CURRENT KNOWLEDGE}

\relprivate{You do not currently know any facts about this problem. You will
have to rely on what other participants tell you.}

\medskip
\textbf{USER MESSAGE}

\textbf{YOUR CURRENT POSITION}

\relprivate{Vote: SOUTHWEST}

\textbf{CURRENT SOCIAL INFORMATION}

\relsocial{Agent 25\\
Vote: B}

\textbf{DECISION}

Work out which option the facts available to you support, and vote for it.

Your reason should briefly explain your choice, for your own record.

Sharing a fact is the only way to pass information to other participants. You
may share exactly one of the facts you currently know by giving its identifier
in \texttt{shared\_fact\_id}, so that the participants who see your position
can use it too. Use ``none'' if you prefer to share nothing.

You may share only a fact listed under YOUR CURRENT KNOWLEDGE, by its exact
identifier. Do not invent facts, identifiers, or relationships that were not
given to you.

Keep your reason to at most three sentences.

Return only valid JSON:

\texttt{\{}\\
\quad\texttt{"vote": "<A | B | C>",}\\
\quad\texttt{"reason": "<brief private reason>",}\\
\quad\texttt{"shared\_fact\_id": "<none>"}\\
\texttt{\}}
\end{relpromptbox}

\paragraph{What changes between calls.}
The identity, query, available semantic answers, private facts, current vote,
visible sources, and admissible fact identifiers are bound from the live game
state. The semantic answers are independently permuted onto A/B/C for every
call. Thus the \relsocial{social vote B} above means \emph{EAST} only in this
instance; the persistent population state stores EAST, not the letter B. This
prevents a stable display position from becoming an unintended convention.

\paragraph{Ordinary and controlled social inputs.}
An uncontrolled update shows the vote and, when supplied, one valid fact from
an ordinary peer. In the audited example the visible source is the controller:
the controller's semantic target EAST is rendered as \relsocial{Vote: B} under
this call's permutation. The prompt deliberately does not identify that source
as a controller or authority. Under the \texttt{recommendation\_only} mode it
also supplies no fact, so the intervention can change the public vote state but
cannot directly enlarge the focal agent's knowledge set.

\paragraph{Initialization variant.}
Before social interaction, each of the 24 agents receives the same role, task,
privacy, fact-sharing, and JSON-contract blocks, but no current-position or
social-information block. Instead, the prompt states that no other participant
has yet declared a position and instructs the agent to decide from its local
facts alone. The resulting 24 ballots define the initial population state.

\paragraph{Response and information boundary.}
The accepted response contains a display-letter vote, a private reason of at
most three sentences, and either one fact identifier already in the agent's
knowledge set or \texttt{none}. The display letter is immediately mapped back
to its semantic relation. The reason is retained for audit but never shown to
another agent. Only the semantic vote and the explicitly selected fact may
enter a later prompt, which makes the public decision channel and the epistemic
transmission channel separately observable.

\paragraph{Audit check.}
The inspection episode recorded 72 distinct compiled prompt instances: 24
local-initialization prompts and 48 focal-update prompts (24 in each of two
population rounds). All 72 responses passed validation without a retry or
provider error. The prompt definition hash was constant across the episode,
whereas every instance hash was unique, confirming that one fixed template was
bound to 72 call-specific views.
